\scanpage{196}

(resp. de rang unipotent nul, resp. affine, resp. unipotent).

\uline{NB} Les cas ``\inserted{connexe} affine'', ``unipotent'' sont triviaux, et essentiellement triviaux \deleted{(le deuxième, \unsure{tout|tous|base} au moins, si l'on \unsure{connait} plus \unsure{générlt|gén.}} (car le premier des deux est \unsure{obvie|évident|achevé|autre}), grâce au th. de Chevalley. Noter que dans le cas unipotent, l'hypothèse signifie que $\omega(s)$ est d'ordre $p^{\unsure{\infty|n}}$ ($p$ la car. en $s$). Le cas \deleted{\unsure{par}} ``rang unipotent nul'' peut se reformuler ainsi, lorsque $E$ est commutatif :

% vertical rule beside the following assertion

Supposons qu'il existe un entier $N$ tel que, pour tt point \uline{fermé} $s$ de $S$, $\omega(s)$ soit \deleted{\unsure{divis.}} dans

\[
(p^NE)(s)=(p^NE_s)(s)=p^NE(s),
\]

i.e. qu'il soit divisible par $p^N$. Alors $\omega(\eta)$ appartient à un sous-groupe algébrique de $E_{\eta}$ de rang \deleted{\unsure{uni}} unipotent nul, donc une puissance de $\omega(\eta)$ appartient à un ss-gpe extension d'une VA par un tore.

% slanted

\margin{Pour ce pb, on peut supposer que $E_\eta$ n'a pas de sous-gpes connexes de rang unipotent nul sauf zéro, \deleted{\unsure{de qq façon|de tte façon}} i.e. \unsure{prend|plus} une extension d'une V.A. par un vectoriel, qui ne splitte sur aucune sous variété abélienne $\ne0$, et il faut conclure que $\omega$, telle que ci-contre, est nécessairement \uline{d'ordre fini}.}

\deleted{Si} ce résultat étant acquis, \unsure{le|ce|les} \deleted{\unsure{question}} cas ``connexe propre'' se ramènerait au cas où $E$ est déjà une extension d'un schéma abélien $A$ par un tore $T$.

\deleted{Si p.ex. $E$ est une extension \uline{non triviale} de \deleted{\unsure{groupe}}\inserted{L'image} \unsure{la} \deleted{\unsure{cas d'une}} \inserted{infini d'} une variété abélienne $A$ par un tore, sur disons, alors il faudrait voir que si $\omega$ est telle que $\omega(s)$} Contrairement aux trois autres cas, ce cas-ci semble également intéressant si on suppose $S$ de car. $p>0$ \unsure{lui-$=$|lui-même} !

% archive sheet number: 195
